\chapter{Pharyngitis}
\index{Pharyngitis}

\medskip
\noindent\textit{Educational information; seek professional advice for individual care.}

\section*{What it is}
Pharyngitis is inflammation of the throat. It is usually caused by a virus, but bacteria such as group A streptococcus can sometimes be responsible. A sore throat may also occur with tonsillitis, glandular fever, allergies, reflux, or irritation from smoke.

\section*{Symptoms}
Symptoms may include a painful or scratchy throat, pain when swallowing, redness, swollen neck glands, fever, tiredness, cough, runny nose, or bad breath. A cough and runny nose make a viral cause more likely, but symptoms alone cannot reliably identify the cause.

\section*{Assessment}
A clinician considers the duration, exposure, fever, rash, cough, swallowing, hydration, immune status, and examination findings. A throat swab or rapid test may be used when bacterial infection is suspected. Antibiotics are not useful for most viral sore throats.

\section*{Self-care and treatment}
Rest, fluids, soft or cool foods, and suitable pain relief can help. Follow the package or clinician’s instructions for medicines, and ask a pharmacist if pregnancy, age, allergies, or other medicines affect what is safe. Confirmed or strongly suspected bacterial infection may need antibiotics.

\section*{When to seek urgent care}
Seek emergency help for difficulty breathing or swallowing, drooling, a high-pitched breathing sound, severe dehydration, confusion, or rapidly worsening illness. Arrange review if symptoms do not improve within about a week, recur often, or occur with a persistent neck lump or mouth ulcer.

\section*{Sources}
\begin{itemize}
\item NHS sore throat guidance: \url{https://www.nhs.uk/symptoms/sore-throat/}
\end{itemize}
